% Part: incompleteness
% Chapter: arithmetization-syntax
% Section: coding-terms

\documentclass[../../../include/open-logic-section]{subfiles}

\begin{document}

\olfileid{inc}{art}{trm}
\olsection{ترميز الحدود}

\begin{explain}
الحدّ ليس إلا نوعًا معيّنًا من متتاليات الرموز: فهو يُبنى استقرائيًا من
الثوابت والمتغيرات وفق قواعد تكوين الحدود. وبما أنه يمكن ترميز متتاليات
الرموز بأعداد---باستخدام مخطط لترميز الرموز وطريقة لترميز متتاليات
الأعداد---فليس من الصعب إسناد أعداد غودل إلى الحدود. وإنما يكمن التحدي
في إثبات أن الخاصية التي يتصف بها عددٌ ما حين يكون عدد غودل لحدّ صحيح
التكوين قابلة للحساب، بل عودية بدائية.
\end{explain}

!!^{variable}s و!!{constant}s أبسط الحدود، ومن السهل اختبار ما إذا كان
$x$ عدد غودل لحد من هذا النوع: تصدق $\fn{Var}(x)$ إذا كان
$x$ هو $\Gn{\Obj v_i}$ لبعض~$i$. وبعبارة أخرى، $x$ متتالية طولها~$1$
وعنصرها الوحيد $(x)_0$ هو شفرة !!{variable} ما~$\Obj v_i$؛ أي إن $x$ يساوي
$\tuple{\tuple{1, i}}$ لبعض~$i$. وبالمثل، تصدق $\fn{Const}(x)$ إذا كان
$x$ هو $\Gn{\Obj c_i}$ لبعض~$i$. وكلتا العلاقتين عودية بدائية، لأنه
إذا وُجد مثل هذا $i$ فلا بد أن يكون $< x$:
\begin{align*}
  \fn{Var}(x) & \defiff \bexists{i<x}{x = \tuple{\tuple{1, i}}}\\
  \fn{Const}(x) & \defiff \bexists{i<x}{x = \tuple{\tuple{2, i}}}
\end{align*}

\begin{prop}
\ollabel{prop:term-primrec}
العلاقتان $\fn{Term}(x)$ و$\fn{ClTerm}(x)$، اللتان تصدقان إذا وفقط إذا
كان $x$ عدد غودل لحد أو لحد مغلق، على الترتيب، عوديتان بدائيتان.
\end{prop}

\begin{proof}
تكون متتالية الرموز~$s$ حدًا إذا وفقط إذا وُجدت متتالية الحدود
$s_0$، \dots، $s_{k-1}=s$ تسجل كيفية تكوين الحد~$s$ من
!!{constant}s و!!{variable}s وفق قواعد تكوين الحدود. ولكي نعبّر عن أن
متتالية تكوين مفترضة كهذه تتبع قواعد التكوين، يلزم أن يتحقق، لكل
$i<k$، أحد الأمور الآتية:
\begin{enumerate}
\item أن يكون $s_i$ !!a{variable}~$\Obj v_j$؛ أو
\item أن يكون $s_i$ !!a{constant}~$\Obj c_j$؛ أو
\item أن يكون $s_i$ مبنيًا من $n$ من الحدود $t_1$، \dots، $t_n$ التي
  ترد قبل الموضع~$i$، باستعمال !!{function}~$\Obj f^n_j$ ذي $n$ مواضع.
\end{enumerate}
ولكي نثبت أن العلاقة المناظرة على أعداد غودل عودية بدائية، علينا أن
نعبّر عن هذا الشرط تعبيرًا عوديًا بدائيًا، أي باستعمال دوال وعلاقات
عودية بدائية وتكميم محدود.

افترض أن $y$ هو العدد الذي يرمّز المتتالية $s_0$، \dots، $s_{k-1}$،
أي إن $y=\tuple{\Gn{s_0}, \dots, \Gn{s_{k-1}}}$. وهو يرمّز متتالية
تكوين للحد ذي عدد غودل~$x$ إذا وفقط إذا تحقق، لكل $i<k$، ما يأتي:
\begin{enumerate}
\item $\fn{Var}((y)_i)$؛ أو
\item $\fn{Const}((y)_i)$؛ أو
\item يوجد $n$ وعدد~$z=\tuple{z_1, \dots, z_n}$ بحيث يساوي كل $z_l$
  قيمةً ما $(y)_{i'}$ من أجل $i'<i$، و
\[
(y)_i = \Gn{\Obj f^n_j(} \concat \fn{flatten}(z) \concat \Gn{)},
\]
\end{enumerate}
وفوق ذلك $(y)_{k-1}=x$. (تحوّل الدالة $\fn{flatten}(z)$ المتتالية
$\tuple{\Gn{t_1}, \dots, \Gn{t_n}}$ إلى $\Gn{t_1, \dots, t_n}$، وهي
عودية بدائية.)

إن المؤشرين $j$ و$n$، وأعداد غودل $z_l$ للحدود $t_l$، وشفرة المتتالية
$z$، أي $\tuple{z_1, \dots, z_n}$، في (3)، كلها أصغر من~$y$. ويمكننا
أن نستبدل $k$ أعلاه بـ$\len{y}$. ومن ثم يمكننا التعبير عن العبارة
«$y$ شفرة متتالية تكوين الحد ذي عدد غودل~$x$» على نحو يبيّن أن هذه
العلاقة عودية بدائية.

ولا يبقى إلا أن نتحقق من وجود حد عودي بدائي لـ$y$. فإذا كان $x$ عدد
غودل لحد، فلا بد أن تكون له متتالية تكوين لا تضم أكثر من $\len{x}$ من
الحدود (إذ يجب أن يبدأ كل حد في متتالية تكوين~$s$ عند موضع ما في~$s$،
ولا يمكن لحدين جزئيين أن يبدآ عند الموضع نفسه). ومن الواضح أن عدد غودل
لكل حد جزئي من~$s$ هو $\le x$. لذلك توجد دائمًا متتالية تكوين شفرتها
$\le p_{k-1}^{k(x+1)}$، حيث $k=\len{x}$.

أما في حالة $\fn{ClTerm}$ فنحذف ببساطة البند الخاص بـ!!{variable}s.
\end{proof}

\begin{prob}
أثبت أن الدالة $\fn{flatten}(z)$، التي تحوّل المتتالية
$\tuple{\Gn{t_1}, \dots, \Gn{t_n}}$ إلى $\Gn{t_1, \dots, t_n}$، عودية
بدائية.
\end{prob}

\begin{prop}\ollabel{prop:num-primrec}
  الدالة $\fn{num}(n)=\Gn{\num{n}}$ عودية بدائية.
\end{prop}

\begin{proof}
  نعرّف $\fn{num}(n)$ بالعودية البدائية:
  \begin{align*}
    \fn{num}(0) & = \Gn{\Obj 0}\\
    \fn{num}(n+1) & = \Gn{\prime(} \concat \fn{num}(n) \concat \Gn{)}.
  \end{align*}
\end{proof}

\end{document}
